% !PRINTED 152
\weberpage{152}
%<<CONT
zeichnung einer Substitution einen einfachen Buchstaben, etwa $A$, und setzt dann

\weq{2}{A = \begin{pmatrix} a_1^{(1)}, & a_2^{(1)}, & \dots & a_n^{(1)} \\ a_1^{(2)}, & a_2^{(2)}, & \dots & a_n^{(2)} \\ \multicolumn{4}{c}{\cdots\cdots\cdots\cdots} \\ a_1^{(n)}, & a_2^{(n)}, & \dots & a_n^{(n)} \end{pmatrix}.}

Wir werden auch abkürzend $A = (a_i^{(\varkappa)})$ setzen und die Determinante der Substitutionscoëfficienten mit $|A|$ bezeichnen. Diese Determinante wird immer von Null verschieden vorausgesetzt. Das Gleichungssystem (1) stellen wir auch symbolisch so dar:

\weq{3}{(y_1, y_2, \dots y_n) = A (x_1, x_2, \dots x_n),}

oder noch kürzer:

\weq{4}{(y) = A (x).}

Die Substitution:

\weq{5}{y_1 = x_1, \quad y_2 = x_2, \quad \dots \quad y_n = x_n}

oder

\weq{6}{J = \begin{pmatrix} 1, & 0, & \dots & 0 \\ 0, & 1, & \dots & 0 \\ \multicolumn{3}{c}{\cdots\cdots\cdots} \\ 0, & 0, & \dots & 1 \end{pmatrix}}

heisst die \gesperrt{identische Substitution}. Eine Substitution der Form

\weq{7}{y_1 = \mu_1 x_1, \quad y_2 = \mu_2 x_2, \quad \dots, \quad y_n = \mu_n x_n}

oder

\weq{8}{M = \begin{pmatrix} \mu_1, & 0, & \dots & 0 \\ 0, & \mu_2, & \dots & 0 \\ \multicolumn{3}{c}{\cdots\cdots\cdots} \\ 0, & 0, & \dots & \mu_n \end{pmatrix}}

soll eine \gesperrt{multiplicative Substitution} oder kurz eine \gesperrt{Multiplication} und die Elemente $\mu_1, \mu_2, \dots, \mu_n$ die \gesperrt{Multiplicatoren} genannt werden, ein Ausdruck, der sich durch die Gleichungen (7) rechtfertigt.

Nehmen wir eine zweite lineare Substitution der Dimension $n$ an, durch die ein neues System von Variablen $z$ eingeführt wird:

\weq{9}{z_i = \sum_{1,n}^{h} b_h^{(i)} y_h, \quad (z) = B (y),}

und führen für $y$ die Werthe aus (1) ein, so erhalten wir die $z$
%CONT>>
